% 张三慧《学习辅导与习题解答》PDF S 独立选答，卷1第4章；未取得V1主教材独立交叉核验。
% 仅录“思考题选答”中实际明确列出的原编号；其余不推断，exam_ready=false。
\section*{第4章\quad 功和能：辅导书已列思考题独立解析}
\paragraph{思考题 4.5}
对于中心力 \(\boldsymbol F=f(r)\hat{\boldsymbol r}\)，沿任何空间微位移的功 \(\boldsymbol F\cdot d\boldsymbol r=f(r)\,dr\)，仅由端点半径决定：\(W=\int_{r_A}^{r_B}f(r)\,dr\)。因此在不穿过奇点且单值的径向力场区域存在 \(\boxed{U(r)=-\int^r f(s)\,ds}\)，力场是保守的。关键是大小\emph{仅依赖r}；非中心的速度依赖力不能仅凭“沿连线”就自动推出保守。

\paragraph{思考题 4.9}
车厢系中悬点静止，摆球所受张力总与相对速度垂直，不做功；若只有重力且阻力忽略，\(\boxed{K'+U_g=\mathrm{const}}\)。地面系中悬点以恒定速度移动，张力可对摆球做功，因 \(\boldsymbol T\cdot(\boldsymbol v'+\boldsymbol V)\) 的第二项不总为零，故只取摆球的 \(K+U_g\) 一般\emph{不守恒}。两惯性系的物理定律一致，但所选择的摆球子系统与能量功项不能混同。

\paragraph{思考题 4.12}
空气动力学的机翼升力来自流体绕翼型的压强分布和动量向下偏转，积分总升力 \(L=\int_S[-p\boldsymbol n+\boldsymbol\tau_{\rm visc}]\cdot\hat z\,dS\)。在特定稳态流线、近似不可压无粘条件下，可用伯努利关系 \(p+\frac12\rho v^2+\rho gz=\mathrm{const}\) 说明较快流速往往对应较低静压，但不能凭“上表面路程较长所以同时到达”推导翼型升力。热气球静浮力则主要因流体静压梯度形成 \(\rho_{\rm fluid}gV\)，机翼是动态压强与动量作用，机理不同。
